\chapter{Mission and Project Management}

\section{Systems Management and Milestones}
The design and launch of a CubeSat 2U platform such as INTISAT is formally articulated under the V-Model of space systems engineering (NASA-SP-2007-6105, ECSS-E-ST-10C). The trajectory follows a strict path: SRR (\textit{System Requirements Review}) → PDR (\textit{Preliminary Design Review}) → \textbf{CDR (Critical Design Review)} of the engineering model → AIT (\textit{Assembly, Integration and Testing}) → FRR (\textit{Flight Readiness Review}) → Flight Model. This document corresponds to the engineering-model CDR.

\section{Schedule, WBS and Coupled Risk Management}
\begin{itemize}
    \item \textbf{Work Breakdown Structure (WBS):} subdivides and parallelizes operational requirements and manufacturing loads into a hierarchical matrix scheme. Branched engineering threads (e.g.\ \textit{OBC software}, \textit{EPS thermal manufacturing}, \textit{AI payload training}) progress encapsulated and converge at transversal mechanical milestones (e.g.\ harness integration and FlatSat tests).

    \item \textbf{Technical Risk Control and Mitigation Matrix:} stochastic foresight tied to hardware failure control. It catalogues potential on-orbit disasters (e.g.\ radiation-induced short-circuits, BGA solder failures after launcher vibration).
\end{itemize}

It also outlines technical pre-flight responses to inevitable contingencies: for example, in case of severe under-voltage due to persistent eclipse, planning delegates by-design preventive intelligence to the \textit{Safe Mode} reaction, stopping fatal degradation, disconnecting collateral instruments and sending emergency requests to the global ground network awaiting troubleshooting on console.
